% !TEX program = xelatex
\documentclass[14pt]{extarticle}

% ============================================================
% GEOMETRY
% ============================================================
\usepackage[
  a4paper,
  left=2.5cm,
  right=1.5cm,
  top=2cm,
  bottom=2cm
]{geometry}

% ============================================================
% МОВА ТА ШРИФТИ
% ============================================================
\usepackage[ukrainian]{babel}
\usepackage{fontspec}
\setmainfont{Times New Roman}

% ============================================================
% ОСНОВНИЙ ТЕКСТ
% ============================================================
\usepackage{setspace}
\onehalfspacing
\setlength{\parindent}{1.27cm}
\setlength{\parskip}{0pt}
\usepackage{microtype}

% ============================================================
% ЗАГОЛОВКИ
% ============================================================
\usepackage{titlesec}

\titleformat{\section}
  {\bfseries\fontsize{16pt}{24pt}\selectfont\centering}
  {\thesection}{0.5em}{}
\titlespacing*{\section}{0pt}{20pt}{6pt}

\titleformat{\subsection}
  {\bfseries\fontsize{14pt}{21pt}\selectfont}
  {\thesubsection}{0.5em}{}
\titlespacing*{\subsection}{1.27cm}{6pt}{6pt}

\titleformat{\subsubsection}
  {\bfseries\itshape\fontsize{14pt}{21pt}\selectfont}
  {\thesubsubsection}{0.5em}{}
\titlespacing*{\subsubsection}{1.27cm}{6pt}{0pt}

% ============================================================
% СПИСКИ
% ============================================================
\usepackage{enumitem}

\setlist[itemize,1]{
  label=$\bullet$,
  leftmargin=1.27cm,
  labelwidth=0.5cm,
  labelsep=0.3cm,
  itemindent=0pt,
  itemsep=3pt,
  topsep=6pt,
  parsep=0pt,
  partopsep=0pt
}

\setlist[enumerate,1]{
  label=\arabic*.,
  leftmargin=2.54cm,
  labelwidth=0.63cm,
  labelsep=0.54cm,
  itemindent=0pt,
  itemsep=0pt,
  topsep=0pt,
  parsep=0pt,
  partopsep=0pt
}

% ============================================================
% ТАБЛИЦІ
% ============================================================
\usepackage{array}
\usepackage{longtable}
\usepackage{booktabs}

% ============================================================
% ГІПЕРПОСИЛАННЯ
% ============================================================
\usepackage[hidelinks,breaklinks=true]{hyperref}

\begin{document}

\sloppy
\emergencystretch=3em

% ==================== ТИТУЛЬНА ЧАСТИНА ====================
\begin{center}
\textbf{Лабораторна робота №1}

\medskip

\textbf{Визначення доцільності створення\\(сфера використання / застосування)}
\end{center}

\medskip

\textbf{Мета роботи:} для планованого проєкту обґрунтувати доцільність створення (сфера використання / застосування) за методами <<Критичних запитань>> та SMART.

\textbf{Проєкт:} <<Система для автоматичної нормалізації суржику та діалектних форм української мови>>.

\textbf{Виконав:} студент групи ШІ-42 Коваль Денис Петрович.

% ==================== РОЗДІЛ 1: КРИТИЧНІ ЗАПИТАННЯ ====================
\section*{1. Метод <<Критичних запитань>>}
\addcontentsline{toc}{section}{1. Метод <<Критичних запитань>>}

Метод <<Критичних запитань>> передбачає послідовну відповідь на ключові запитання щодо проєкту, які дозволяють визначити його доцільність, межі та зацікавлені сторони.

\subsection*{1.1. Яку проблему вирішує проєкт?}

Українська мова характеризується високою варіативністю: поряд із літературною нормою поширені регіональні діалекти (гуцульський, бойківський, закарпатський тощо) та російсько-український суржик. Внутрішня міграція, яка посилилась після початку повномасштабної війни, призводить до того, що люди з різних регіонів частіше опиняються в нових мовних середовищах, де діалектні слова, суржикові конструкції або локальні форми можуть бути незрозумілими у побуті, навчанні чи роботі.

Водночас наявні інструменти автоматичної корекції української мови орієнтовані переважно на стандартні орфографічні та граматичні помилки, тоді як семантична нормалізація діалектних або суржикових конструкцій підтримується рідко і часто реалізується правиловими (rule-based) підходами, які не забезпечують достатнього покриття.

Таким чином, проєкт вирішує проблему \textbf{відсутності автоматизованого інструменту для нормалізації діалектних та суржикових форм до літературної української мови} з використанням сучасних нейромережевих підходів.

\subsection*{1.2. Які потреби задовольняє проєкт?}

Проєкт задовольняє такі потреби:
\begin{itemize}
  \item потреба у зрозумілому спілкуванні між носіями різних діалектів та літературної мови --- абітурієнти, студенти, внутрішні мігранти потребують простого інструменту для нормалізації коротких текстів (повідомлень, оголошень, чатів);
  \item потреба розробників NLP-систем у модулі нормалізації --- діалектні та суржикові вхідні тексти знижують якість роботи пайплайнів, побудованих на літературній українській;
  \item потреба у підвищенні якості україномовних текстових даних --- нормалізація покращує якість корпусів для подальшого навчання моделей;
  \item потреба дослідників у відтворюваному baseline-рішенні для задачі діалектної нормалізації української мови.
\end{itemize}

\subsection*{1.3. Яка сфера застосування проєкту?}

Сфера застосування проєкту охоплює:
\begin{itemize}
  \item нормалізацію коротких текстів (повідомлення, оголошення, чати, коментарі) з діалектних форм та суржику до літературної української;
  \item підтримку декількох українських діалектів: гуцульського (як базового), а також бойківського та закарпатського;
  \item підтримку нормалізації російсько-українського суржику;
  \item надання веб-інтерфейсу для вибору типу вхідного тексту (конкретний діалект або суржик) та отримання нормалізованого результату;
  \item використання як окремого інструменту або як компоненту більших NLP-систем.
\end{itemize}

\subsection*{1.4. Хто є зацікавленими сторонами?}

\begin{itemize}
  \item \textbf{Кінцеві користувачі} --- студенти, абітурієнти та мігранти з діалектних регіонів, яким потрібен інструмент для нормалізації повсякденних текстів;
  \item \textbf{Дослідники NLP} --- науковці, які працюють із обробкою природної мови для української, зокрема з низькоресурсними мовними варіантами;
  \item \textbf{Розробники програмного забезпечення} --- інженери, яким потрібен модуль нормалізації як компонент у ширшій системі (чат-боти, пошукові системи, системи аналізу тексту);
  \item \textbf{Науковий керівник та кафедра} --- академічна зацікавленість у результатах дослідження та апробації методів.
\end{itemize}

\subsection*{1.5. Що є за межами сфери застосування?}

За межами проєкту залишаються:
\begin{itemize}
  \item повний переклад з інших мов (російської, польської тощо) --- система працює лише з діалектами та суржиком;
  \item нормалізація довгих документів (наукових статей, книг) --- орієнтація на короткі тексти;
  \item розпізнавання усного мовлення (ASR) --- система працює виключно з текстовими вхідними даними;
  \item повне комерційне розгортання з масштабуванням на великі навантаження;
  \item покриття всіх діалектів української мови --- лише обрані діалекти та суржик;
  \item створення вичерпного лінгвістичного корпусу --- використовуються наявні ресурси та синтетичне розширення.
\end{itemize}

\subsection*{1.6. Трикутник конкурентних обмежень}

Трикутник обмежень (час, ресурси, якість) для проєкту:

\begin{itemize}
  \item \textbf{Час} --- проєкт обмежений термінами виконання бакалаврської кваліфікаційної роботи (один навчальний семестр). Етапи: побудова baseline-моделі для гуцульського діалекту, масштабування на мультидіалектний корпус, донавчання LLM, створення інтерфейсу.

  \item \textbf{Ресурси} --- обчислювальні ресурси обмежені: baseline-модель (seq2seq) тренується на CPU; донавчання Lapa LLM (12B параметрів) потребує GPU з $\geq$24~GB відеопам'яті та техніки QLoRA для зменшення вимог до пам'яті. Дані обмежені: для більшості діалектів відсутні великі ручні корпуси, тому використовується синтетичне розширення.

  \item \textbf{Якість} --- цільова якість вимірюється метриками BLEU, chrF++, TER та COMET. Baseline на гуцульському діалекті досягає BLEU~$\approx$51. Контрольні набори формуються з ручних (не синтетичних) прикладів для забезпечення об'єктивності оцінки.
\end{itemize}

\textbf{Пріоритетність:} у контексті бакалаврської роботи пріоритет надається якості дослідження та відтворюваності результатів, потім --- дотриманню термінів, і останнім --- мінімізації ресурсних витрат (за рахунок QLoRA замість повного донавчання).

% ==================== РОЗДІЛ 2: SMART ====================
\section*{2. Метод SMART}
\addcontentsline{toc}{section}{2. Метод SMART}

Метод SMART дозволяє оцінити мету проєкту за п'ятьма критеріями: конкретність (Specific), вимірюваність (Measurable), наявність виконавця (Assignable), реалістичність (Realistic), обмеженість у часі (Time-related).

\subsection*{2.1. S --- Specific (Конкретність)}

Мета проєкту сформульована конкретно: \textit{розробити та експериментально оцінити систему автоматичної нормалізації українських діалектів (гуцульського, бойківського, закарпатського) та російсько-українського суржику до літературної української мови на основі нейронних моделей}.

Конкретність забезпечується чітким визначенням:
\begin{itemize}
  \item вхідних даних --- тексти діалектами та суржиком;
  \item вихідних даних --- нормалізований текст літературною українською;
  \item підходу --- seq2seq baseline + інструкційне донавчання Lapa LLM із LoRA/QLoRA;
  \item результату --- працююча система з веб-інтерфейсом.
\end{itemize}

\subsection*{2.2. M --- Measurable (Вимірюваність)}

Результати проєкту є вимірюваними за допомогою стандартних метрик оцінювання якості машинного перекладу та текстового редагування:
\begin{itemize}
  \item \textbf{BLEU} --- модифікована n-грамна точність; baseline для гуцульського діалекту досягає $\approx$51;
  \item \textbf{chrF++} --- символьна F-міра, більш придатна для морфологічно багатих мов;
  \item \textbf{TER} --- кількість редагувань, необхідних для приведення гіпотези до еталону;
  \item \textbf{COMET} --- нейромережева метрика, що корелює з людськими оцінками якості.
\end{itemize}

Контрольні набори (validation/test) сформовані з вручну анотованих пар, що забезпечує об'єктивність вимірювань.

\subsection*{2.3. A --- Assignable (Наявність виконавця)}

Проєкт виконується студентом групи ШІ-42 Ковалем Денисом Петровичем у рамках бакалаврської кваліфікаційної роботи за спеціальністю 122 <<Комп'ютерні науки>> під керівництвом наукового керівника. Виконавець має необхідні компетенції у галузі машинного навчання, обробки природної мови та програмування мовою Python.

Розподіл відповідальності:
\begin{itemize}
  \item \textbf{Виконавець} --- збір і підготовка даних, реалізація моделей, проведення експериментів, створення інтерфейсу, оформлення результатів;
  \item \textbf{Науковий керівник} --- методичне спрямування, контроль якості дослідження, рецензування.
\end{itemize}

\subsection*{2.4. R --- Realistic (Реалістичність)}

Реалістичність проєкту обґрунтовується такими факторами:
\begin{itemize}
  \item \textbf{Наявність даних} --- для гуцульського діалекту доступні вручну анотований паралельний корпус, синтетичний корпус та діалектний словник; для інших діалектів і суржику використовується синтетичне розширення на базі літературних речень із датасету Spivavtor;
  \item \textbf{Перевірені методи} --- підхід базується на усталених техніках: seq2seq Transformer для baseline, інструкційне донавчання LLM, LoRA/QLoRA для параметро-ефективного навчання;
  \item \textbf{Наявність базової моделі} --- Lapa LLM є відкритою україномовною моделлю з адаптованим токенізатором, що суттєво спрощує донавчання для української;
  \item \textbf{Достатні обчислювальні ресурси} --- baseline тренується на CPU, а QLoRA дозволяє донавчати 12B-модель на GPU з 24~GB відеопам'яті;
  \item \textbf{Методологічна підтримка} --- існують релевантні дослідження (Vuyko Mistral для діалектного перекладу, UA-GEC для граматичної корекції), які підтверджують працездатність обраного підходу.
\end{itemize}

\subsection*{2.5. T --- Time-related (Обмеженість у часі)}

Проєкт обмежений термінами виконання бакалаврської кваліфікаційної роботи. Етапи виконання:
\begin{enumerate}
  \item Аналіз предметної області та огляд літератури.
  \item Збір та підготовка даних для гуцульського діалекту; побудова та оцінка baseline-моделі.
  \item Формування мультидіалектного корпусу; синтетичне розширення для бойківського, закарпатського діалектів та суржику.
  \item Інструкційне донавчання Lapa LLM з LoRA/QLoRA; експериментальне оцінювання.
  \item Розробка веб-інтерфейсу (Gradio) та інтеграція системи.
  \item Оформлення результатів та підготовка до захисту.
\end{enumerate}

% ==================== ВИСНОВКИ ====================
\section*{Висновки}
\addcontentsline{toc}{section}{Висновки}

У ході виконання лабораторної роботи було обґрунтовано доцільність створення системи для автоматичної нормалізації суржику та діалектних форм української мови за двома методами.

\textbf{Метод <<Критичних запитань>>} дозволив визначити: проблему (відсутність інструментів нормалізації діалектів і суржику), потреби користувачів (зрозуміле спілкування, компонент NLP-пайплайнів), сферу застосування (короткі тексти, декілька діалектів, веб-інтерфейс), зацікавлені сторони (кінцеві користувачі, дослідники, розробники) та межі проєкту. Аналіз трикутника обмежень показав, що пріоритет надається якості дослідження з дотриманням академічних термінів.

\textbf{Метод SMART} підтвердив, що мета проєкту є конкретною (чітко визначені вхід/вихід, підхід, результат), вимірюваною (BLEU, chrF++, TER, COMET), забезпеченою виконавцем (студент під керівництвом наукового керівника), реалістичною (наявність даних, перевірених методів, базових моделей) та обмеженою у часі (семестр виконання бакалаврської роботи).

Обидва методи свідчать про доцільність реалізації проєкту.

\end{document}
